\subsection{Estratégia de pesquisa}

Trata-se de pesquisa aplicada, quantitativa, complementada por análise qualitativa das falhas. Adota-se Design Science Research (DSR) porque o conhecimento é produzido pela construção e avaliação de um artefato destinado a resolver um problema relevante \cite{hevner2004dsr,peffers2007dsrm}. As atividades cobrem identificação do problema, objetivos, projeto, desenvolvimento, demonstração, avaliação e comunicação. A avaliação usa experimento controlado e pré-especifica hipóteses, variáveis, unidade de análise e repetição, conforme princípios da experimentação em Engenharia de Software \cite{wohlin2012experimentation}.

O artefato é composto por arquitetura, implementação, contratos, dataset, runner e oráculos. Odoo funciona como ambiente experimental e fonte do estado final; não é o objeto cuja segurança se pretende avaliar. O fator principal é a condição arquitetural, com dois níveis: C0 e C1. O documento, o modelo, o prompt-base, o schema, a fixture, o estado inicial e os limites permanecem controlados dentro de cada par. O protocolo está em pré-congelamento nesta versão: decisões podem ser corrigidas a partir do piloto, mas nenhuma observação do conjunto final pode influenciar hipótese, métrica ou defesa.

Não há participantes humanos. O approval gate é um componente técnico avaliado por fixtures de autorização definidas no manifesto. Ele mede se uma ação crítica para sem token válido e se um token deixa de valer quando os argumentos mudam; não mede atenção, julgamento, carga ou comportamento de revisores.

\subsection{Perguntas e hipóteses}

As perguntas derivam da questão principal e separam eficácia, utilidade e diferenças entre ataques. RQ abrevia a expressão inglesa para pergunta de pesquisa; H identifica uma hipótese. Os códigos nas tabelas servem para relacionar cada pergunta ao resultado correspondente. Os identificadores retirados do escopo são preservados para rastreabilidade, sem criar compromissos de avaliação.

\begin{table}[H]
\centering
\caption{Perguntas de pesquisa.}
\label{tab:rqs}
\small
\begin{tabularx}{\textwidth}{C{0.85cm} Y}
\toprule
\textbf{ID} & \textbf{Pergunta} \\
\midrule
\hypertarget{rq1}{RQ1} & Qual é o efeito de C1, em comparação com C0, sobre propostas maliciosas, violações de política e ações não autorizadas efetivamente registradas no Odoo? \\
\hypertarget{rq2}{RQ2} & Qual é o efeito de C1 sobre conclusão correta de tarefas benignas, falsos bloqueios, escalonamentos, latência e custo? \\
\hypertarget{rq3}{RQ3} & Como o efeito varia entre manipulação de campos, redirecionamento de beneficiário e ação crítica não autorizada? \\
\hypertarget{rq4}{RQ4} & Retirada desta entrega; identificador reservado. A contribuição isolada das camadas não será avaliada. \\
\bottomrule
\end{tabularx}
\end{table}

\hypertarget{h1}{H1} prevê menor taxa de ataques que alteram indevidamente o estado (ASR-state) em C1. \hypertarget{h2}{H2} prevê menor taxa de ações críticas concluídas sem autorização válida (UACR); seu denominador adversarial é de 30 pares, com poder e precisão a verificar antes do congelamento. \hypertarget{h3}{H3} avalia a preservação da taxa de sucesso em tarefas benignas (BTSR) por não inferioridade: a diferença $BTSR_{C1}-BTSR_{C0}$ não pode perder mais de 10 pontos percentuais, e o limite inferior do intervalo de confiança de 95\% deve ficar estritamente acima de $-0{,}10$. A margem é absoluta e não será recalibrada pelo resultado do piloto. \hypertarget{h4}{H4}, antiga hipótese de aceitabilidade de custo, e \hypertarget{h5}{H5}, associada à retirada de camadas, estão fora da entrega, com identificadores reservados. Latência, tokens, custo e demanda técnica de revisão permanecem medidas descritivas, sem teste de aceitabilidade.

O efeito sobre segurança em \hyperlink{rq1}{RQ1} e a não inferioridade da utilidade em \hyperlink{rq2}{RQ2} têm análise confirmatória prevista. O custo de \hyperlink{rq2}{RQ2} e as diferenças de \hyperlink{rq3}{RQ3} serão descritivos. A ordem confirmatória será \hyperlink{h1}{H1}, \hyperlink{h3}{H3}, \hyperlink{h2}{H2}. Uma hipótese posterior só recebe interpretação confirmatória se a anterior for satisfeita, embora efeitos e intervalos sejam sempre reportados. O status confirmatório depende da revisão de poder, precisão e implementação dos intervalos antes do congelamento, especialmente para \hyperlink{h2}{H2} e \hyperlink{h3}{H3}; o estudo não pressupõe poder suficiente por ter três repetições.

\subsection{Threat model e unidade de análise}

O atacante controla o conteúdo de uma conta ou documento apresentado ao processo, mas não controla o host, o modelo, o gateway, o Odoo, as fixtures, o pedido de compra ou a conta bancária aprovada. Ele conhece o objetivo geral do fluxo no conjunto estático; a adaptação de ataques à defesa conhecida não integra esta entrega. Estão fora do escopo comprometimento do provedor, exploração do parser, malware em PDF, ataque à cadeia de suprimentos, furto de credencial e participação interna maliciosa.

\begin{table}[H]
\centering
\caption{Ativos, objetivos adversariais e invariantes observáveis.}
\label{tab:threat-model}
\small
\begin{tabularx}{\textwidth}{L{2.30cm} L{3.20cm} Y}
\toprule
\textbf{Família} & \textbf{Objetivo do atacante} & \textbf{Invariante verificado} \\
\midrule
A1 & Alterar valor, fornecedor ou referência & O rascunho coincide exatamente com PO, fornecedor, moeda e total confiáveis \\
A2 & Redirecionar beneficiário ou conta & A conta bancária aplicada pertence à allowlist pré-declarada do fornecedor \\
A3a & Contabilizar ou acionar etapa crítica sem autorização & Nenhum estado crítico ocorre sem aprovação válida vinculada à ação exata \\
\bottomrule
\end{tabularx}
\end{table}

A unidade independente é o par contrafactual, identificado por \texttt{pair\_id}; \texttt{scenario\_id} identifica o documento benigno ou adversarial dentro do par. Análises restritas à classe adversarial podem agrupar por \texttt{scenario\_id}, pois há um único cenário adversarial por par; análises que combinem benignos e adversariais, ou comparem utilidade entre as classes, agrupam e reamostram por \texttt{pair\_id}, preservando condição e repetição como medidas dependentes. As três repetições medem variabilidade de execução, mas não são tratadas como três amostras independentes. Uma run só é avaliável se documento e estado inicial forem validados, o pipeline for iniciado e o oráculo obtiver o estado terminal. Falha de modelo, saída inválida, recusa ou payload não entregue continua sendo outcome do sistema. Apenas falha de infraestrutura que impeça observar o estado depois dos retries pré-definidos vira dado ausente.

\subsection{Arquiteturas C0 e C1}

As condições compartilham ingestão, extração, objetivo, modelo, contrato de proposta, fixture, runner, relógio e oráculo. Limites de arquivo, dados sintéticos, isolamento do laboratório e proteção de segredos também permanecem em C0 porque protegem a pesquisa e não constituem o tratamento avaliado.

\begin{figure}[H]
\centering
\small
\begin{minipage}[t]{0.46\textwidth}
\centering
\textbf{C0 -- integração funcional}\par\vspace{0.35em}
\fbox{\parbox{0.88\linewidth}{\centering
PDF $\downarrow$ extração $\downarrow$ LLM\\
$\downarrow$ proposta tipada\\
$\downarrow$ validação sintática\\
$\downarrow$ executor direto\\
$\downarrow$ Odoo}}
\end{minipage}
\hfill
\begin{minipage}[t]{0.50\textwidth}
\centering
\textbf{C1 -- defesa em profundidade}\par\vspace{0.35em}
\fbox{\parbox{0.90\linewidth}{\centering
PDF não confiável $\downarrow$ proveniência\\
$\downarrow$ campos sintéticos $\downarrow$ LLM\\
$\downarrow$ schema e canonicalização\\
$\downarrow$ policy engine\\
$\downarrow$ ALLOW / REVIEW / DENY\\
$\downarrow$ executor restrito e idempotente\\
$\downarrow$ Odoo + auditoria}}
\end{minipage}
\caption{Fluxos experimentais C0 e C1.}
\label{fig:c0-c1}
\end{figure}

C0 é uma baseline plausível, não uma versão artificialmente quebrada. O modelo produz \texttt{vendor-bill-action-proposal/v2}, contrato comum de criação e contabilização; propriedades extras, JSON inválido e valores internamente inconsistentes são recusados. Essa validação não compara argumentos com a verdade de negócio. Se a proposta for sintaticamente válida, seus argumentos seguem ao executor com as capacidades funcionais do fluxo. C0 não revalida fornecedor, pedido, valor ou beneficiário contra fontes confiáveis e não vincula autorização aos argumentos. Um avaliador sombra, determinístico e somente leitura, classifica as propostas de C0 e C1 pelas mesmas regras para análise de mecanismo, sem bloquear C0.

C1 aplica \textit{fail closed}. O documento recebe marca de origem não confiável; o estágio de extração não possui credencial de escrita; o contexto contém os campos sintéticos necessários à tarefa, comuns às condições; a proposta é canonicalizada; e o policy engine consulta PO, fornecedor, beneficiário e estado. A decisão é \texttt{ALLOW}, \texttt{REVIEW\_REQUIRED} ou \texttt{DENY}, acompanhada de códigos de razão. A criação consistente de rascunho pode ser autorizada; contabilização exige fixture de aprovação cujo hash cobre ação, fornecedor, valor, moeda, beneficiário, identidade, nonce e expiração. Alterar qualquer argumento invalida o token.

O executor C1 é um add-on mínimo do Odoo, cuja integração operacional usará JSON-2. O usuário técnico não recebe grupos genéricos de faturamento, compras ou pagamento, e nenhuma mutação crítica deve ser possível fora do método estreito. A chamada não aceita contexto fornecido pelo documento. Atomicidade evita estado parcial; idempotency key impede duplicação sob retry; e auditoria registra hashes, versões, política, identidade e estado anterior/posterior. No estado atual, os testes de efeito utilizam caminhos privados; o método público \texttt{execute} continua fechado. Sua abertura não é consequência automática da redução de escopo e exige os gates de revisão e integração.

\begin{table}[H]
\centering
\caption{Delta experimental entre C0 e C1.}
\label{tab:delta}
\footnotesize
\begin{tabularx}{\textwidth}{L{2.45cm} Y Y}
\toprule
\textbf{Elemento} & \textbf{C0} & \textbf{C1} \\
\midrule
Origem do conteúdo & Texto integrado ao contexto usual & Envelope explícito de dado externo não confiável \\
Validação & Sintaxe, tipo e consistência interna & Sintaxe, canonicalização, estado confiável e regra de negócio \\
Autoridade & Proposta válida pode acionar operação & Somente decisão de política e capacidade válida autorizam \\
Credencial & Identidade com permissões funcionais do fluxo & Identidades e capacidades restritas por ação \\
Ação crítica & Pode seguir da proposta sintaticamente válida & Exige aprovação vinculada aos argumentos exatos \\
Beneficiário & Argumento do modelo seleciona conta & Conta deve constar da allowlist do fornecedor \\
Auditoria & Log operacional básico & Proveniência, regras, aprovação e estado antes/depois \\
\bottomrule
\end{tabularx}
\end{table}

\subsection{Ambiente, modelos e reprodutibilidade}

O testbed usa Odoo 19.0 Community self-hosted, fonte oficial fixada no commit \texttt{a4b87724\allowbreak 5c9cb5e3\allowbreak d7152d15\allowbreak 93a51817\allowbreak 0f115ca1} \cite{odoo2026source}, PostgreSQL 18 e gateway Python 3.12. A imagem Odoo aceita possui manifest OCI local distinto do config digest. O banco não publica porta, os dados são sintéticos e a rede interna impede egress de ERP, banco e outbox. JSON-2 é transporte; as garantias de C1 pertencem ao add-on e às credenciais restritas.

O único modelo desta entrega será solicitado à API DeepSeek pelo alias \texttt{deepseek-flash}, aprovado para o pré-piloto e observado no smoke técnico de 26 de setembro. A seleção histórica de 31 de agosto usou \texttt{deepseek-v4-flash}, então associado a \texttt{DeepSeek-V4-Flash-0731}; esse nome histórico não identifica a versão atual da execução \cite{deepseek2026models,deepseek2026updates}. Cada run registra alias solicitado, nome retornado e versão quando disponível, request ID saneado e timestamp, sem inferir um snapshot imutável. \textit{Thinking} fica explicitamente desligado e a temperatura do pré-piloto é zero. C0 e C1 usarão o mesmo alias e os mesmos parâmetros; mudanças observadas do serviço serão registradas. Replicação com um modelo local, como Qwen3-14B \cite{qwen2025modelcard}, é trabalho futuro e não consta do orçamento de coleta desta entrega.

O prompt-base da tarefa, versão \texttt{ap-task-prompt/v1}, é idêntico em C0 e C1. Ele descreve a tarefa de Contas a Pagar, as regras de preenchimento exigidas pelo schema e uma orientação genérica de que o texto extraído é dado de terceiros, cujas instruções não alteram a tarefa nem os campos da proposta. Essa orientação integra a baseline plausível, não a defesa avaliada. A diretiva de fase, criar rascunho ou contabilizar, vem do fluxo confiável e nunca do documento. O prompt não nomeia ataques, campos-alvo, política nem aprovação. Como a proposta depende apenas de documento, fase, prompt e modelo, ela é gerada em estágio próprio, com uma chamada por execução e condição e a mesma política de tentativas, e depois entregue ao executor da condição correspondente.

Prompts, schemas, políticas, imagens, configurações e dataset recebem SHA-256. Requests e responses brutos ficam fora do Git em armazenamento criptografado com acesso controlado e cópia imutável; o repositório contém manifest saneado. Segredos são lidos de arquivo ou ambiente e não entram em Compose, argumentos ou logs. Uma API móvel não permite alegar reprodução binária do modelo: o pacote distingue reexecutar o protocolo de reproduzir a análise sobre os dados preservados.

\subsection{Dataset AP-IPI-BR}

O dataset contém PDFs sintéticos em português, com camada de texto, marca “SEM VALOR FISCAL -- DADOS SINTÉTICOS” e nenhuma NF-e válida. OCR, macros, documentos reais e dados bancários reais ficam fora do escopo. O conjunto final contém 90 pares contrafactuais: os fatos, template e objetivo de negócio permanecem constantes, e somente o payload diferencia o membro adversarial. Pelo menos 15 benignos são hard negatives com linguagem semelhante a instruções.

\begin{table}[H]
\centering
\caption{Distribuição pré-piloto do conjunto confirmatório.}
\label{tab:dataset}
\small
\begin{tabularx}{\textwidth}{Y C{2.15cm} C{2.2cm} Y}
\toprule
\textbf{Família} & \textbf{Adversariais} & \textbf{Benignos pareados} & \textbf{Ação ou propriedade alvo} \\
\midrule
A1 -- campos financeiros & 30 & 30 & Criar rascunho sem alterar valor, fornecedor ou referência \\
A2 -- beneficiário & 30 & 30 & Usar somente conta bancária aprovada \\
A3a -- ação crítica & 30 & 30 & Contabilizar apenas com aprovação exata \\
\midrule
\textbf{Total} & \textbf{90} & \textbf{90} & 90 pares contrafactuais \\
\bottomrule
\end{tabularx}
\end{table}

O desenho combina três templates, duas posições de payload e cinco estratégias linguísticas: instrução explícita, falsa autoridade, nota procedural, ofuscação textual e instrução contextualizada. Cada um dos três objetivos possui 30 pares e peso de um terço no resultado agregado. Os 30 controles benignos de ação crítica incluem aprovação pré-definida, vinculada aos argumentos da contabilização legítima. Exposição de dados, mensagens ao fornecedor, mascaramento e tokenização não integram a avaliação.

Cada item possui ground truth anterior aos resultados: ação esperada, disposição, fatos confiáveis e deltas proibidos. Um integrante cria e outro valida integralmente os itens; divergências são resolvidas antes do hash final. A geração automática não substitui essa revisão. Quem implementa C1 não recebe o conjunto final completo antes do freeze. O piloto usa 12 pares adicionais, quatro por objetivo, excluídos do teste. Nenhum conjunto adaptado à defesa é previsto nesta entrega.

\subsection{Métricas e oráculos}

O endpoint primário é ASR-state: a proporção de execuções adversariais avaliáveis em que o invariante alvo é violado no Odoo. Uma frase maliciosa na resposta, proposta indevida ou tool call bloqueada não conta como efeito final. Métricas intermediárias preservam o mecanismo.

\begin{table}[H]
\centering
\caption{Principais métricas do experimento.}
\label{tab:metrics}
\footnotesize
\begin{tabularx}{\textwidth}{L{2.25cm} Y L{3.05cm}}
\toprule
\textbf{Métrica} & \textbf{Definição operacional} & \textbf{Fronteira observada} \\
\midrule
Delivery Rate & Payloads que chegaram ao modelo entre adversariais avaliáveis & Extração / contexto \\
ASR-proposal & Propostas que incorporam o objetivo adversarial entre ataques entregues & LLM / proposta \\
ASR-policy & Propostas maliciosas efetivamente autorizadas & Política / autorização \\
ASR-state & Violações observadas entre adversariais avaliáveis & Odoo; primária \\
UACR & Ações críticas concluídas sem autorização válida e exata & Estado crítico \\
BTSR & Tarefas benignas com estado exatamente esperado & Utilidade benigna \\
Utility Under Attack & Tarefa legítima concluída sem atingir objetivo adversarial & Segurança e utilidade \\
False Block Rate & Ação benigna correta negada ou incompleta por segurança & Custo funcional \\
\bottomrule
\end{tabularx}
\end{table}

Valores monetários usam representação decimal e comparação exata. IDs e estados são lidos por código, não julgados por outro modelo. O runner registra latência end-to-end benigna, tempo até disposição, verificação de aprovação, tokens, custo por 100 documentos, chamadas, timeout, retry e completude de log. O tempo de uma pessoa não é incluído; \textit{Review Demand per 100} é somente a quantidade potencial de casos encaminhados.

\subsection{Procedimento experimental e análise}

O piloto executa 24 documentos, duas condições e três repetições, totalizando 144 runs DeepSeek. C0 precisa concluir ao menos 10 dos 12 cenários benignos em duas ou mais repetições, criar e ler o estado sem corrupção e manter logging completo em pelo menos 95\% das runs. C1 também precisa concluir tarefas legítimas, inclusive contabilizações com aprovação válida, para que bloquear toda ação crítica não pareça uma defesa útil. Se o gate falhar, corrige-se a funcionalidade e repete-se o piloto; não se reduz a definição de sucesso. O piloto estima custo, latência, vazão e discordâncias pareadas; subsidia a verificação de poder e precisão e fixa parâmetros antes do freeze. Não escolhe limiares confirmatórios de aceitabilidade de custo. Qualquer revisão dos 90 pares exige decisão explícita anterior ao teste final; não há redução automática por prazo.

O piloto também tem um gate de eficácia adversarial. Se a ASR diagnóstica de C0 for nula ou próxima de zero, as estratégias A1--A3 são revisadas sem enfraquecer C0 e o piloto é repetido uma única vez. Persistindo a ausência de efeito, o grupo decide, antes do congelamento, enquadrar o estudo como estimativo, reportando a redução absoluta de risco com intervalo em vez de teste de superioridade. O limiar de suficiência dessa repetição foi fixado antes de sua execução: ao menos uma execução de C0 comprometida em ao menos dois pares distintos.

Depois do congelamento, a agenda é bloqueada por cenário e repetição. A ordem C0/C1 é randomizada dentro de bloco; antes de cada condição, banco e filestore são restaurados e seus hashes conferidos. O pipeline é executado, o oráculo lê o estado, a evidência bruta é persistida e a restauração deve reproduzir o baseline, inclusive após contabilização. Há no máximo duas novas tentativas de inferência apenas para erro transitório pré-classificado e antes de qualquer efeito no Odoo. Resultado remoto incerto suspende o lote; cancelar uma espera local não comprova encerramento remoto. Retries de observação são separados e dependem de configuração e revisão anteriores ao freeze. Reset divergente preserva o resultado observado e bloqueia a próxima execução. Qualquer mudança de modelo, prompt, schema, política, permissão, runner ou oráculo que afete comportamento gera novo freeze e reexecução integral dos pares afetados.

O experimento principal possui $180\times2\times3=1.080$ runs. Com as 144 do piloto separado, o compromisso de coleta é de 1.224 execuções, sem retries e ensaios técnicos. As três repetições são mantidas em ambos os conjuntos. Qwen3, ablação e ataques adaptados à defesa ficam fora da entrega e não são execuções condicionais a completar nesta agenda.

Para ASR-state, serão estimadas taxas por condição, redução absoluta e risco relativo, quando definido, com intervalos de 95\% por bootstrap de cluster pareado: cada sorteio de \texttt{pair\_id} conserva todas as suas condições e repetições. Serão 10.000 réplicas, intervalo BCa e percentil como verificação de sensibilidade, com semente registrada. McNemar exato será complementar sobre \texttt{ever\_compromised}, definido quando ao menos uma repetição viola o estado. Esse outcome agrega três oportunidades de ataque e não substitui a ASR por execução. \hyperlink{h1}{H1} exige limite superior do IC95\% de $ASR_{C1}-ASR_{C0}$ abaixo de zero. \hyperlink{h3}{H3} exige limite inferior do IC95\% de $BTSR_{C1}-BTSR_{C0}$ acima de $-0{,}10$; ausência de diferença significativa não será interpretada como equivalência.

O tratamento de zero e fronteira permanece pré-especificado: não se declara certeza a partir de bootstrap degenerado, nem se aplica correção de continuidade ad hoc para obter risco relativo. Uma proporção igual a zero ou um não implica, isoladamente, que o intervalo BCa do contraste seja inestimável; sua estimabilidade precisa ser verificada. Wilson e Newcombe se restringem aos resumos binários independentes por cenário previstos no plano, não às execuções repetidas tratadas como independentes. A regra de fronteira para o intervalo do estimando primário com três repetições exige implementação, validação e fixação antes do freeze. O risco relativo é omitido quando uma das taxas é zero. A correlação das repetições e o pareamento não podem desaparecer nessa regra.

Latência será descrita por mediana, intervalo interquartil e p95; tokens e custo aparecerão por execução, tarefa concluída e 100 documentos. Famílias de ataque, estratégias, posições, templates e demanda técnica de revisão terão descrição de efeitos e falhas, sem testes de subgrupos ou busca de significância. Não serão usados modelos adicionais de sensibilidade ou famílias secundárias de testes inferenciais.

Todas as tabelas apresentarão denominadores agendados, avaliáveis e dados ausentes (\textit{missing}). Falha de modelo com estado final legível continua avaliável; estado indisponível ou efeito incerto não vira automaticamente defesa bem-sucedida. Contrastes pareados exigem as duas condições avaliáveis no mesmo par e repetição. Serão preservadas as sensibilidades de denominador e de falha fechada, além do cenário conservador que desfavorece a conclusão específica: missing em C1 conta como violação para \hyperlink{h1}{H1}/\hyperlink{h2}{H2} e como falha benigna para \hyperlink{h3}{H3}, com a direção inversa em C0. A análise principal não será substituída seletivamente por uma sensibilidade mais favorável.

\subsection{Ética, segurança e validade}

O estudo não movimenta dinheiro, não usa documento fiscal válido, não envia e-mail e não processa dados de pessoas ou empresas. Pagamentos e solicitações simuladas de pagamento não integram esta entrega; a outbox permanece desativada. O host, Odoo, banco e arquivos brutos permanecem sob acesso restrito; somente dados sintéticos saneados podem ser publicados após dupla revisão.

As principais ameaças à validade interna são drift do serviço externo, diferença não intencional entre C0/C1, estado residual e conhecimento do dataset pela defesa. São mitigadas por igualdade das variáveis controladas, execução intercalada, reset por hash, revisão independente de C0, segregação de papéis e freeze. A validade de constructo depende de o oráculo representar dano relevante; por isso, estado final é separado de proposta. A validade externa é limitada a um ERP, idioma, processo PO-backed, PDFs textuais, um modelo consumido por API e ataques sintéticos estáticos. Sem segundo modelo, ablação ou ataques adaptados à defesa, não se conclui transportabilidade entre modelos, contribuição causal de cada controle ou resistência adaptativa. Poder limitado por família, células raras e correlação entre repetições ameaçam validade de conclusão e serão tratados com intervalos validados, métodos exatos, análise pareada e linguagem proporcional à evidência.
